% Single-file MDPI / Applied Sciences manuscript.
% Use this file inside the official MDPI Overleaf template project.
% It avoids external table/bib dependencies and removes unsupported metadata commands.

\documentclass[applsci,article,submit,pdftex,moreauthors]{Definitions/mdpi}

\firstpage{1}
\makeatletter
\setcounter{page}{\@firstpage}
\makeatother
\pubvolume{16}
\issuenum{1}
\articlenumber{0}
\pubyear{2026}
\copyrightyear{2026}
\externaleditor{Academic Editor: TBD}
\datereceived{}
\daterevised{}
\dateaccepted{}
\datepublished{}

\usepackage{booktabs}
\usepackage{threeparttable}
\usepackage{multirow}
\usepackage{makecell}
\usepackage{adjustbox}
\usepackage{array}
\usepackage{amsmath}
\usepackage{amssymb}
\usepackage{graphicx}
\usepackage{enumitem}
\usepackage{microtype}
\usepackage{float}

\Title{Exploratory Artificial Intelligence for Automated Social Media Advertising Content Creation and Multi-Objective Creative Optimization}

\Author{Author One $^{1}$, Author Two $^{1,*}$ and Author Three $^{2}$}
\AuthorNames{Author One, Author Two and Author Three}

\address{%
$^{1}$ \quad School/Department, University Name, City, Postcode, Country; author1@email.com (A.O.); author2@email.com (A.T.)\\
$^{2}$ \quad School/Department, University Name, City, Postcode, Country; author3@email.com (A.T.)}

\corres{Correspondence: corresponding.author@email.com}

\abstract{Generative artificial intelligence has substantially accelerated digital-media content production, yet many operational workflows still rely on single-pass generation followed by manual screening. That pattern increases throughput but does not provide a principled way to explore creative alternatives, compare them under shared objectives, or reproduce selection decisions. This study proposes \textbf{EAI-CO} (Exploratory AI-based Creative Optimization), a framework that formulates social-media advertising content creation as an exploratory multi-objective optimization problem rather than a one-shot generation task. EAI-CO encodes campaign briefs into structured task representations, generates diverse creative candidates, evaluates them through a compliance-aware multi-objective scoring function, and iteratively refines high-potential candidates through controlled search. The main benchmark is conducted on a reproducible local \texttt{Qwen2.5-7B-Instruct} backbone over 10 products, 4 audience segments, and 40 product--audience tasks. EAI-CO is compared against template-based, single-shot, open-source-only, and prompt-engineered baselines, together with four mechanistic ablations. EAI-CO achieves the best overall reward (0.8051) and significantly outperforms all primary baselines ($p < 0.001$ to $p < 10^{-12}$). The full method also outperforms all ablated variants, indicating that diversity preservation, iterative refinement, and audience modeling each contribute materially to final quality. A supplementary cross-model validation on \texttt{gpt-5.4-mini-2026-03-17} preserves the same ranking pattern and supports the transferability of the framework to a stronger commercial generator. Overall, the main gain in digital-media ad generation does not arise only from stronger backbones; it arises from treating creative production as an exploratory optimization process with explicit control over audience fit, diversity, and compliance-aware quality.}

\keyword{exploratory artificial intelligence; generative AI; digital media; social media advertising; automated content creation; creative optimization}

\begin{document}

\section{Introduction}

Digital-media advertising increasingly depends on the rapid production of personalized creative assets across platforms, products, and audience segments. Social-media advertisements, in particular, must communicate clearly under severe space constraints while remaining visually coherent, audience-relevant, and distinguishable from competing messages. This creates a structural tension for marketing operations: advertisers need both scale and relevance, whereas manual creative development remains expensive, slow, and difficult to standardize.

Generative artificial intelligence offers a direct response to this bottleneck. It can produce headlines, captions, visual prompts, and promotional variants at a speed that conventional workflows cannot match. However, speed alone is not the central scientific issue. In many practical deployments, generative models are still used in a narrow one-shot manner: a prompt is written, one or a few candidates are produced, and human operators manually select among them. Although this increases throughput, it does not solve the deeper optimization problem. A single-pass workflow neither systematically explores the creative search space nor explains why one candidate should be preferred to another. As a result, evaluation remains ad hoc and reproducibility remains weak.

This paper argues that automated advertising content creation should be treated as an exploratory optimization problem rather than as a direct generation problem. In actual campaign practice, creatives are rarely judged in isolation. Marketers compare variants that differ in emotional framing, benefit emphasis, call-to-action strategy, and audience alignment. A rigorous AI system should therefore do more than generate fluent copy. It should generate alternative candidates, preserve meaningful diversity, evaluate them under explicit objectives, and iteratively refine the most promising directions.

To address this need, we propose \textbf{EAI-CO} (Exploratory AI-based Creative Optimization), a framework for automated social-media advertising content creation and optimization. EAI-CO integrates structured campaign-brief encoding, exploratory candidate generation, automatic multi-objective evaluation, iterative refinement, and reproducible experiment management. The framework is designed to move AI-assisted advertising away from isolated prompt--response interactions and toward a controlled search-and-selection process.

The present study makes four contributions. First, it reformulates social-media ad generation as an exploratory multi-objective optimization problem. Second, it develops a compliance-aware evaluator that jointly considers relevance, clarity, aesthetic plausibility, audience fit, diversity, safety, and predicted engagement. Third, it provides a reproducible benchmark protocol with baseline comparisons, mechanistic ablations, audience-wise analyses, and runtime reporting. Fourth, it validates transferability by showing that the same superiority pattern holds under both a local open-source backbone and a stronger commercial generator.

\section{Related Work}

\subsection{Generative AI in Marketing and Advertising}

Generative AI is reshaping marketing by expanding the feasible scale of content production, personalization, and campaign experimentation. Existing marketing scholarship has emphasized both the opportunities and the governance challenges of this transition. Generative models can accelerate copywriting, ideation, and promotional adaptation, but marketing applications require stronger controls than general-purpose text generation. In particular, marketing creatives must remain aligned with brand positioning, audience expectations, factual constraints, and measurable business goals \cite{grewal2025,heitmann2024}.

This distinction matters because a generative system that is merely fluent is not necessarily useful in advertising. Content must also be audience-sensitive, operationally differentiable, and sufficiently safe for deployment. The present work builds on this observation by treating advertising generation as a constrained optimization problem rather than as unconstrained text synthesis.

\subsection{AI-Generated Advertising Creatives}

Recent work on AI-generated advertising creatives spans creative generation, creative ranking, ad-creative selection, and click-through-rate-oriented optimization. Industrial systems increasingly treat creatives as decision variables that interact with larger ad-serving pipelines. Prior studies have examined the joint optimization of ad ranking and creative selection, as well as parallel ranking architectures for ads and creatives in real-time systems \cite{lin2022,yang2024aaai}. These studies demonstrate that creative choice is not peripheral to advertising performance; it is a core component of it.

At the same time, much of this literature depends on proprietary datasets, industrial infrastructure, or online feedback that is unavailable to academic replication. This makes it difficult to compare creative-generation methods under controlled and reproducible conditions. The present paper contributes a framework-oriented benchmark that remains tractable under local execution while preserving explicit links to practical advertising objectives.

\subsection{Prompt Optimization and Iterative Generation}

Prompt engineering has become a standard practical technique for improving large-model outputs, including marketing copy. However, prompt optimization is often under-specified as a research methodology. It may improve results without making the search process itself observable. Iterative generation partially addresses this issue by revising prompts or outputs across multiple rounds, but many iterative systems still lack an explicit account of diversity preservation and candidate selection.

Our approach extends this line of work by embedding iterative generation within a structured search space. Candidates are not merely regenerated; they are explored under defined creative axes and evaluated through a shared reward function. This makes it possible to analyze what is gained by search rather than by prompt fluency alone.

\subsection{Multi-Objective Evaluation for Creative Systems}

Advertising quality is intrinsically multi-objective. A useful creative should be relevant to the product, clear to the viewer, visually plausible, aligned with the target audience, distinct from alternative variants, and compliant with factual or safety constraints. A single scalar criterion such as fluency, lexical diversity, or model confidence is therefore insufficient. Our framework follows the multi-objective view and explicitly models creative quality as the result of balancing competing objectives rather than maximizing a single generation score.

\section{Materials and Methods}

\subsection{Problem Definition}

Let a campaign task be defined by a product $p$, an audience segment $a$, a platform context $s$, and a constraint set $c$. The system must generate a social-media advertisement
\begin{equation}
x = \{\text{headline}, \text{body}, \text{cta}, \text{visual\_prompt}\},
\end{equation}
with the aim of maximizing a multi-objective reward while remaining aligned with the campaign brief. The key shift is that the objective is not to generate a single plausible creative, but to identify the strongest creative from an explored candidate set.

\subsection{Campaign Brief Encoder}

The campaign-brief encoder converts each task into a structured schema containing product title, category, selling points, audience segment, platform, tone, and constraints. This representation standardizes the task description across methods and ensures that differences among outputs arise from model behavior and search strategy rather than inconsistent prompt framing.

For example, a portable espresso maker can be paired with the \texttt{young\_professionals} segment and the platform context \texttt{Instagram-style social feed}, whereas constraint fields can encode safety requirements such as avoiding unsupported health or medical claims. This structure supports audience-aware generation while preserving fairness across methods.

\subsection{Exploratory Candidate Generator}

Instead of querying the model for one direct answer, EAI-CO samples a controlled creative space. Each candidate is defined by a set of exploration axes, including emotional style, appeal type, layout style, color direction, CTA style, caption length, and audience pain point.

For each product--audience task, the framework generates multiple candidates by varying these axes. Under the final benchmark configuration, two rounds are used with two candidates per round, and one elite candidate is retained for refinement. This design yields a shallow but explicit exploratory loop that is computationally manageable while still exposing the effect of search.

Each candidate contains four fields: headline, body, CTA, and visual prompt. Although the current benchmark does not execute a full large-scale image-generation experiment, the visual-prompt field is preserved because multimodal coherence remains part of the creative specification.

\subsection{Compliance-Aware Multi-Objective Evaluator}

Each candidate is scored on seven positive dimensions and one penalty term: relevance, clarity, aesthetic quality, audience fit, predicted engagement, diversity, brand safety, and factuality penalty.

For the final local benchmark, the reward weights are as follows: relevance $=0.20$, clarity $=0.16$, aesthetic $=0.14$, audience fit $=0.16$, predicted engagement $=0.16$, diversity $=0.08$, and brand safety $=0.10$.

The reward is computed as
\begin{equation}
R(x) = 0.20r_{\mathrm{rel}} + 0.16r_{\mathrm{clar}} + 0.14r_{\mathrm{aes}} + 0.16r_{\mathrm{aud}} + 0.16r_{\mathrm{eng}} + 0.08r_{\mathrm{div}} + 0.10r_{\mathrm{safe}} - p_{\mathrm{fact}}.
\end{equation}

This formulation is intentionally compliance-aware. It does not reward persuasive intensity in isolation. Instead, it balances persuasive potential with audience alignment and plausible, safe wording.

\subsection{Iterative Optimization Loop}

EAI-CO proceeds in three conceptual stages:
\begin{enumerate}[leftmargin=1.5em]
    \item \textbf{Exploration}: diverse initial candidates are sampled from the creative space.
    \item \textbf{Refinement}: high-scoring elites are retained and mutated under controlled variation.
    \item \textbf{Selection}: the resulting candidates are scored and the strongest creative is selected under the shared reward.
\end{enumerate}

The framework therefore behaves as an optimizer rather than as a simple conditional generator. It assumes that strong creatives emerge from structured comparison and refinement, not from accepting the first fluent answer.

\subsection{Baselines and Ablations}

We compare EAI-CO against four primary baselines: \texttt{B0\_Template}, \texttt{B1\_SingleShot\_API}, \texttt{B2\_OpenSource\_Only}, and \texttt{B3\_PromptEngineered\_AI}. To isolate the contribution of internal mechanisms, we also evaluate four ablations: \texttt{Ours\_without\_diversity}, \texttt{Ours\_without\_factual\_penalty}, \texttt{Ours\_without\_iterative\_loop}, and \texttt{Ours\_without\_audience\_modeling}.

\section{Experimental Setup}

\subsection{Benchmark Construction}

The main benchmark comprises 10 products crossed with 4 audience segments: students, young professionals, family users, and price-sensitive consumers. This yields 40 product--audience tasks. Under the local main experiment, 9 final method outputs are retained per task, corresponding to 5 primary methods and 4 ablations, for a total of 360 evaluated creatives.

\subsection{Model Settings}

\subsubsection{Main Reproducible Backbone}

The main benchmark uses a local \texttt{Qwen2.5-7B-Instruct} backbone executed on an RTX 3090 GPU. This choice prioritizes reproducibility, local control, and scientific traceability.

\subsubsection{Commercial Transfer Backbone}

A supplementary cross-model validation is conducted on \texttt{gpt-5.4-mini-2026-03-17}. This experiment preserves the same 10 products and 4 audience segments but evaluates only the 5 primary methods. Its purpose is to test framework transferability rather than to replace the main benchmark.

\subsection{Search Configuration}

For both the local benchmark and the commercial transfer experiment, the exploration settings are fixed to 2 rounds, 2 candidates per round, and 1 elite candidate retained for refinement. This configuration provides measurable exploratory behavior without introducing impractically high computational cost.

\subsection{Evaluation Protocol}

All reported metrics are derived from the shared evaluator defined in Section 3.4. Statistical comparisons are conducted between \texttt{Ours\_EAI\_CO} and each primary baseline over the matched task set. The reported outputs include mean reward, reward deltas, significance values, predicted-engagement deltas, audience-fit deltas, audience-wise summaries, and runtime statistics.

\subsection{Reproducibility and Verification}

The project includes fixed configuration files, executable benchmark scripts, generated CSV summaries, and per-method analysis scripts. Large outputs are excluded from version control for repository hygiene, but the experimental workflow itself is script-driven and reproducible.

\subsection{Figure Design Notes}

Figure~\ref{fig:eai_co_pipeline} summarizes the EAI-CO pipeline from campaign-brief encoding to exploratory candidate generation, multi-objective evaluation, iterative refinement, and final selection. Figure~\ref{fig:candidate_evolution} illustrates candidate evolution under the iterative optimization loop for a representative product--audience task.

\begin{figure}[H]
\centering
\fbox{\rule{0pt}{2.0in}\rule{0.92\linewidth}{0pt}}
\caption{Overview of the EAI-CO framework. The pipeline begins with campaign-brief encoding, proceeds through exploratory candidate generation and compliance-aware multi-objective evaluation, and ends with iterative refinement and final selection.}
\label{fig:eai_co_pipeline}
\end{figure}

\begin{figure}[H]
\centering
\fbox{\rule{0pt}{2.0in}\rule{0.92\linewidth}{0pt}}
\caption{Example candidate evolution under the iterative optimization loop. The figure should illustrate how early-round variants differ from the final selected creative under the same product--audience task.}
\label{fig:candidate_evolution}
\end{figure}

\begin{table*}[t]
\centering
\caption{Benchmark configuration and compared methods.}
\label{tab:benchmark_setup}
\begin{threeparttable}
\begin{tabular}{p{2.8cm} p{11.8cm}}
\toprule
Component & Specification \\
\midrule
Main backbone & Local \texttt{Qwen2.5-7B-Instruct} executed on an RTX 3090 GPU \\
Cross-model validation backbone & Commercial \texttt{gpt-5.4-mini-2026-03-17} \\
Products & 10 \\
Audience segments & Students, young professionals, family users, and price-sensitive consumers \\
Tasks & 40 product-audience tasks \\
Primary baselines & \texttt{B0\_Template}, \texttt{B1\_SingleShot\_API}, \texttt{B2\_OpenSource\_Only}, \texttt{B3\_PromptEngineered\_AI}, and \texttt{Ours\_EAI\_CO} \\
Ablations & \texttt{Ours\_without\_diversity}, \texttt{Ours\_without\_factual\_penalty}, \texttt{Ours\_without\_iterative\_loop}, and \texttt{Ours\_without\_audience\_modeling} \\
Search depth & 2 rounds, 2 candidates per round, and 1 elite candidate retained for refinement \\
Reward weights (main benchmark) & relevance $0.20$, clarity $0.16$, aesthetic $0.14$, audience fit $0.16$, predicted engagement $0.16$, diversity $0.08$, and brand safety $0.10$ \\
Main evaluated outputs & 360 creatives for the local main experiment; 200 creatives for the commercial cross-model validation \\
\bottomrule
\end{tabular}
\end{threeparttable}
\end{table*}

\begin{table*}[t]
\centering
\caption{Main benchmark results on the local \texttt{Qwen2.5-7B-Instruct} backbone. Best values among the primary methods are shown in bold.}
\label{tab:qwen_main_results}
\begin{adjustbox}{max width=\textwidth}
\begin{tabular}{lccccccccc}
\toprule
Method & Reward & Relevance & Clarity & Audience Fit & Diversity & Pred.\ Eng. & Brand Safety & Penalty & Latency (ms) \\
\midrule
\texttt{Ours\_EAI\_CO} & \textbf{0.8051} & \textbf{0.9122} & 0.9130 & \textbf{0.4719} & 0.7693 & \textbf{0.8238} & \textbf{0.9945} & 0.0010 & 11114.23 \\
\texttt{B3\_PromptEngineered\_AI} & 0.7568 & 0.7920 & \textbf{0.9165} & 0.4394 & 0.8600 & 0.7849 & 0.9395 & 0.0110 & 2727.48 \\
\texttt{B1\_SingleShot\_API} & 0.7494 & 0.7920 & 0.9130 & 0.4367 & \textbf{0.8660} & 0.7695 & 0.9175 & 0.0130 & 2780.30 \\
\texttt{B2\_OpenSource\_Only} & 0.7445 & 0.8050 & 0.8955 & 0.4312 & 0.8631 & 0.7646 & 0.9340 & 0.0120 & 2707.25 \\
\texttt{B0\_Template} & 0.6856 & 0.5580 & 0.9200 & 0.4042 & 1.0000 & 0.6570 & 1.0000 & 0.0000 & 1.00 \\
\bottomrule
\end{tabular}
\end{adjustbox}
\end{table*}

\begin{table*}[t]
\centering
\caption{Ablation results on the local \texttt{Qwen2.5-7B-Instruct} benchmark. Negative reward gaps indicate degradation relative to the full EAI-CO framework.}
\label{tab:qwen_ablation}
\begin{adjustbox}{max width=\textwidth}
\begin{tabular}{lcccccc}
\toprule
Method & Reward & Reward Gap vs. Ours & Audience Fit & Audience-Fit Gap & Pred.\ Eng. & Brand Safety \\
\midrule
\texttt{Ours\_EAI\_CO} & \textbf{0.8051} & 0.0000 & \textbf{0.4719} & 0.0000 & \textbf{0.8238} & 0.9945 \\
\texttt{w/o factual penalty} & 0.8006 & -0.0045 & 0.4692 & -0.0027 & 0.8185 & 0.8130 \\
\texttt{w/o audience modeling} & 0.7833 & -0.0218 & 0.3960 & -0.0759 & 0.8100 & \textbf{1.0000} \\
\texttt{w/o iterative loop} & 0.7760 & -0.0291 & 0.4421 & -0.0298 & 0.7959 & 0.9670 \\
\texttt{w/o diversity} & 0.7349 & -0.0702 & 0.4556 & -0.0163 & 0.8149 & 0.9890 \\
\bottomrule
\end{tabular}
\end{adjustbox}
\end{table*}

\begin{table*}[t]
\centering
\caption{Audience-wise reward on the local \texttt{Qwen2.5-7B-Instruct} benchmark. Best values within each audience segment are shown in bold.}
\label{tab:qwen_audience}
\begin{adjustbox}{max width=\textwidth}
\begin{tabular}{lcccc}
\toprule
Method & Students & Young Professionals & Family Users & Price-sensitive Consumers \\
\midrule
\texttt{Ours\_EAI\_CO} & \textbf{0.8098} & \textbf{0.8050} & \textbf{0.8147} & \textbf{0.7908} \\
\texttt{B3\_PromptEngineered\_AI} & 0.7169 & 0.7554 & 0.7785 & 0.7764 \\
\texttt{B1\_SingleShot\_API} & 0.7648 & 0.7359 & 0.7389 & 0.7580 \\
\texttt{B2\_OpenSource\_Only} & 0.7229 & 0.7385 & 0.7663 & 0.7503 \\
\texttt{B0\_Template} & 0.6797 & 0.6875 & 0.6992 & 0.6758 \\
\bottomrule
\end{tabular}
\end{adjustbox}
\end{table*}

\begin{table*}[t]
\centering
\caption{Cross-model validation results on \texttt{gpt-5.4-mini-2026-03-17}. Best values among the primary methods are shown in bold.}
\label{tab:gpt_cross_model}
\begin{adjustbox}{max width=\textwidth}
\begin{tabular}{lccccccccc}
\toprule
Method & Reward & Relevance & Clarity & Audience Fit & Diversity & Pred.\ Eng. & Brand Safety & Penalty & Latency (ms) \\
\midrule
\texttt{Ours\_EAI\_CO} & \textbf{0.9957} & \textbf{1.0000} & \textbf{0.7390} & \textbf{0.7562} & 0.8041 & \textbf{0.8506} & \textbf{1.0000} & 0.0000 & 17128.73 \\
\texttt{B2\_OpenSource\_Only} & 0.9751 & 0.9870 & 0.7365 & 0.6452 & 0.8631 & 0.8145 & \textbf{1.0000} & 0.0000 & 6443.85 \\
\texttt{B1\_SingleShot\_API} & 0.9736 & 0.9870 & 0.6935 & 0.7183 & \textbf{0.8660} & 0.8211 & 0.9890 & 0.0110 & 6188.20 \\
\texttt{B3\_PromptEngineered\_AI} & 0.9729 & 0.9838 & 0.6925 & 0.6994 & 0.8600 & 0.8276 & 0.9890 & 0.0120 & 10776.03 \\
\texttt{B0\_Template} & 0.8396 & 0.5580 & 0.9200 & 0.4042 & 1.0000 & 0.6570 & \textbf{1.0000} & 0.0000 & 1.00 \\
\bottomrule
\end{tabular}
\end{adjustbox}
\end{table*}

\begin{table*}[t]
\centering
\caption{Statistical comparison of EAI-CO against primary baselines and runtime profile. Left: paired significance tests on the local \texttt{Qwen2.5-7B-Instruct} benchmark. Right: average runtime and model-call profile.}
\label{tab:significance_cost}
\begin{minipage}[t]{0.60\textwidth}
\centering
\begin{threeparttable}
\begin{tabular}{lccc}
\toprule
Baseline & $\Delta$ Reward & $p$-value & $\Delta$ Audience Fit \\
\midrule
\texttt{B0\_Template} & +0.1195 & $9.09\times10^{-13}$ & +0.0677 \\
\texttt{B1\_SingleShot\_API} & +0.0557 & $1.34\times10^{-6}$ & +0.0352 \\
\texttt{B2\_OpenSource\_Only} & +0.0606 & $5.18\times10^{-8}$ & +0.0406 \\
\texttt{B3\_PromptEngineered\_AI} & +0.0483 & $6.58\times10^{-5}$ & +0.0325 \\
\bottomrule
\end{tabular}
\end{threeparttable}
\end{minipage}
\hfill
\begin{minipage}[t]{0.36\textwidth}
\centering
\begin{threeparttable}
\begin{tabular}{lcc}
\toprule
Method & Calls & Mean Latency (ms) \\
\midrule
\texttt{Ours\_EAI\_CO} & 8.0 & 11114.23 \\
\texttt{B3\_PromptEngineered\_AI} & 4.0 & 2727.48 \\
\texttt{B1\_SingleShot\_API} & 4.0 & 2780.30 \\
\texttt{B2\_OpenSource\_Only} & 3.0 & 2707.25 \\
\texttt{B0\_Template} & 0.0 & 1.00 \\
\bottomrule
\end{tabular}
\end{threeparttable}
\end{minipage}
\end{table*}

\section{Results}

\subsection{Main Benchmark Results on the Local Backbone}

Table~\ref{tab:qwen_main_results} reports the main benchmark results obtained on the local \texttt{Qwen2.5-7B-Instruct} backbone. EAI-CO achieved the highest overall reward (\textbf{0.8051}), outperforming \texttt{B3\_PromptEngineered\_AI} (0.7568), \texttt{B1\_SingleShot\_API} (0.7494), \texttt{B2\_OpenSource\_Only} (0.7445), and \texttt{B0\_Template} (0.6856). As summarized in Table~\ref{tab:significance_cost}, the reward gains over all primary baselines were statistically significant. Relative to \texttt{B0\_Template}, EAI-CO improved reward by +0.1195 ($p = 9.09 \times 10^{-13}$). The gains over the stronger single-pass baselines also remained significant: +0.0557 over \texttt{B1\_SingleShot\_API} ($p = 1.34 \times 10^{-6}$), +0.0606 over \texttt{B2\_OpenSource\_Only} ($p = 5.18 \times 10^{-8}$), and +0.0483 over \texttt{B3\_PromptEngineered\_AI} ($p = 6.58 \times 10^{-5}$).

Beyond the scalar reward, the full framework achieved the best audience-fit score (0.4719) and the strongest brand-safety score (0.9945) among the primary methods. This indicates that the observed advantage is not restricted to a single dimension such as relevance or engagement, but reflects a more balanced optimization outcome.

\subsection{Mechanistic Evidence from the Ablation Study}

Table~\ref{tab:qwen_ablation} presents the ablation study. The full method outperformed all ablated variants. The largest degradation occurred when diversity preservation was removed (-0.0702), indicating that search-space breadth is a major contributor to the final gain. Removing the iterative loop reduced reward by -0.0291, showing that controlled refinement contributes beyond one-pass generation. Removing audience modeling lowered reward by -0.0218 and reduced audience fit from 0.4719 to 0.3960, directly supporting the claim that persona conditioning is functionally important.

The factual-penalty ablation remained close to the full model (-0.0045). This result should be interpreted cautiously. It does not imply that compliance-aware controls are unnecessary. Rather, it indicates that factual safeguards can be integrated with only a small tradeoff in persuasive reward while preserving a more disciplined optimization process.

\subsection{Audience-Wise Robustness}

Table~\ref{tab:qwen_audience} reports the audience-wise reward values on the local benchmark. EAI-CO ranked first in all four audience groups: family users (0.8147), students (0.8098), young professionals (0.8050), and price-sensitive consumers (0.7908). This consistency reduces the concern that the overall effect is driven by a single favorable subgroup and suggests that the exploratory optimization strategy remains effective across materially different audience profiles.

\subsection{Runtime Characteristics}

Table~\ref{tab:significance_cost} summarizes runtime behavior. EAI-CO required 8.0 mean model calls per task and 11114.23 ms mean latency, whereas the single-pass baselines required approximately 3--4 model calls and 2700--2800 ms mean latency. This overhead is expected because the framework explicitly trades additional inference budget for exploratory search quality. For the intended offline campaign-design setting, this tradeoff remains operationally acceptable.

\subsection{Cross-Model Validation on \texttt{gpt-5.4-mini-2026-03-17}}

Table~\ref{tab:gpt_cross_model} presents the fixed-subset cross-model validation on \texttt{gpt-5.4-mini-2026-03-17}. The same ranking pattern observed on the local backbone was preserved: EAI-CO achieved the highest reward (\textbf{0.9957}), followed by \texttt{B2\_OpenSource\_Only} (0.9751), \texttt{B1\_SingleShot\_API} (0.9736), \texttt{B3\_PromptEngineered\_AI} (0.9729), and \texttt{B0\_Template} (0.8396). The gains over all baselines remained statistically significant, with reward improvements of +0.1561 ($p = 1.54 \times 10^{-8}$) over \texttt{B0\_Template}, +0.0221 ($p = 1.41 \times 10^{-3}$) over \texttt{B1\_SingleShot\_API}, +0.0206 ($p = 6.42 \times 10^{-4}$) over \texttt{B2\_OpenSource\_Only}, and +0.0228 ($p = 1.01 \times 10^{-3}$) over \texttt{B3\_PromptEngineered\_AI}.

The audience-wise rankings were also preserved, with EAI-CO remaining first for family users (0.9950), price-sensitive consumers (0.9938), students (1.0000), and young professionals (0.9939). These results support the claim that the advantage of the proposed exploratory optimization framework transfers to a stronger commercial model under the same task protocol.

\section{Discussion}

\subsection{The Primary Gain Comes from Optimization Structure}

The most stable pattern across the local benchmark and the commercial-model validation is that exploratory optimization outperforms one-shot generation. This result indicates that the central contribution of EAI-CO lies less in any specific generator and more in the optimization structure itself. In digital-media advertising, creative quality should therefore be treated as a search-and-selection problem rather than as a direct text-generation problem.

\subsection{Diversity and Iteration Are Material Design Factors}

The ablation study shows that diversity preservation and iterative refinement are the main structural sources of improvement. This aligns with practical creative work: strong campaign creatives rarely emerge from the first acceptable draft; they emerge from comparing strategically different variants and refining the strongest direction. The present benchmark makes this process measurable and reproducible.

\subsection{Audience Modeling Has Functional Value}

The audience-modeling ablation provides direct empirical evidence that personalization is not simply a rhetorical layer. Once persona conditioning is removed, both reward and audience fit decline. This implies that the framework is not merely producing generic persuasive copy but is measurably exploiting audience-specific information to alter the optimization trajectory.

\subsection{Compliance-Aware Scoring Is Practically Necessary}

The factual-penalty ablation remained close to the full framework, indicating that compliance-aware optimization can be introduced without substantially damaging persuasive performance. This is important for practical deployment. In real digital-media settings, a method that yields slightly stronger promotional language at the cost of weaker factual discipline would not necessarily be preferable.

\subsection{Cross-Model Transfer Strengthens the Contribution}

The commercial-model experiment should be interpreted as a transferability test rather than as the main source of effect-size contrast. Because \texttt{gpt-5.4-mini} is a stronger generator under the current evaluator, reward values become compressed near the upper bound. Nevertheless, preservation of the full ranking pattern materially strengthens the paper by reducing the risk that the main result is an artifact of one local open-source backbone.

\subsection{Limitations}

This study has several limitations. First, it does not use real platform click-through-rate or conversion outcomes. The claims should therefore remain limited to automatic reward superiority, predicted engagement, audience-conditioned quality, and cross-model ranking consistency. Second, the current submission package does not include a completed human-evaluation study. This does not invalidate the automatic and transfer results, but it does narrow the claims that can be made. Third, the benchmark is methodologically useful but remains smaller than a production-scale industrial evaluation environment. Future work should expand product diversity, multimodal richness, and deployment realism.

\section{Conclusions}

This paper presented EAI-CO, a framework that treats social-media advertising content creation as an exploratory multi-objective optimization problem. Across a reproducible local benchmark on \texttt{Qwen2.5-7B-Instruct}, the framework significantly outperformed template-based and single-pass baselines and remained superior under a commercial cross-model validation on \texttt{gpt-5.4-mini-2026-03-17}. The ablation study showed that diversity preservation, iterative refinement, and audience modeling each contribute materially to the observed gains.

The broader implication is that the main scientific contribution of AI-assisted advertising may lie less in scaling model size alone and more in designing explicit optimization frameworks that search, evaluate, and refine creative alternatives under measurable constraints. This provides a more rigorous foundation for future work on digital-media content generation, audience-aware creative design, and responsible automated marketing systems.

\authorcontributions{Conceptualization, A.O. and A.T.; methodology, A.O. and A.T.; software, A.O.; validation, A.O. and A.T.; formal analysis, A.O.; investigation, A.O.; resources, A.T.; data curation, A.O.; writing---original draft preparation, A.O.; writing---review and editing, A.O., A.T. and A.T.; visualization, A.O.; supervision, A.T.; project administration, A.T. All authors have read and agreed to the published version of the manuscript.}

\funding{This research received no external funding.}

\institutionalreviewboardstatement{Not applicable.}

\informedconsentstatement{Not applicable.}

\dataavailabilitystatement{The code, configuration files, and structured benchmark summaries supporting the reported results are available in the project repository. Large generated outputs are omitted from version control for repository hygiene but can be made available by the corresponding author upon reasonable request.}

\acknowledgments{The authors acknowledge the use of local GPU resources for the reproducible open-source benchmark and API-based access for the commercial-model transfer experiment.}

\conflictsofinterest{The authors declare no conflict of interest.}

\reftitle{References}
\begin{thebibliography}{999}

\bibitem{grewal2025}
Grewal, D.; Satornino, C.B.; Davenport, T.; Guha, A.
How Generative AI Is Shaping the Future of Marketing.
\emph{J. Acad. Mark. Sci.} \textbf{2025}, \emph{53}, 1220--1242.

\bibitem{heitmann2024}
Heitmann, M.
Generative AI for Marketing Content Creation: New Rules for an Old Game.
\emph{NIM Mark. Intell. Rev.} \textbf{2024}, \emph{16}, 10--17.

\bibitem{lin2022}
Lin, K.; Zhang, X.; Li, F.; Wang, P.; Long, Q.; Deng, H.; Xu, J.; Zheng, B.
Joint Optimization of Ad Ranking and Creative Selection.
In \emph{Proceedings of the 45th International ACM SIGIR Conference on Research and Development in Information Retrieval}; ACM: New York, NY, USA, 2022; pp. 2627--2631.

\bibitem{yang2024aaai}
Yang, Z.; Sang, L.; Wang, H.; Chen, W.; Wang, L.; He, J.; Peng, C.; Lin, Z.; Gan, C.; Shao, J.
Parallel Ranking of Ads and Creatives in Real-Time Advertising Systems.
In \emph{Proceedings of the AAAI Conference on Artificial Intelligence}; AAAI Press: Washington, DC, USA, 2024.

\bibitem{yang2024www}
Yang, H.; Yuan, J.; Yang, S.; Xu, L.; Yuan, S.; Zeng, Y.
A New Creative Generation Pipeline for Click-Through Rate with Stable Diffusion Model.
In \emph{Companion Proceedings of the ACM Web Conference 2024}; ACM: New York, NY, USA, 2024.

\end{thebibliography}

\end{document}
